% !TeX program = lualatex
% Vietnamese translation for OI Public Library
% Source-PDF: IOI中国国家候选队论文/2000/陈彧论文.pdf
\documentclass[11pt]{article}
\usepackage{oipl-vn}

\title{Phương pháp tư duy trong thi tin học}
\author{Chen Yu}
\date{}

\begin{document}
\maketitle

\origtitle{Thinking methods in informatics competitions}
\sourcepath{IOI China candidate team papers / 2000 / Chen Yu.pdf}

\paragraph{Từ khóa.}
Tin học; phương pháp tư duy.

\section*{Tóm tắt}

Bài viết này mượn một số lý thuyết về tư duy toán học để bàn về vị trí và tác
dụng của phương pháp tư duy trong thi tin học. Nội dung chính giới thiệu vài
phương pháp tư duy thường gặp trong thi tin học, gồm: thử nghiệm, phỏng đoán
và quy nạp; mô hình hóa; phân loại và chia để trị; tương tự. Bài viết dùng
nhiều ví dụ để phân tích quá trình suy nghĩ, phần lớn lấy từ các đề NOI và
IOI năm 1999, nên có tính đại diện khá rộng. Cuối cùng là phần tổng kết mục
đích và gợi mở của bài viết.

\section{Dẫn nhập}

``Olympic là thể dục của tư duy''. Cũng như các kỳ thi học thuật khác, thi tin
học, trên nền tảng tri thức phong phú và năng lực thực hành nhất định, là cuộc
đọ sức của tư duy với tư duy. Những thí sinh xuất sắc thường có suy nghĩ nhanh
và linh hoạt. Vì vậy, trong khi liên tục thảo luận phương pháp và tổng kết
kinh nghiệm, ta cũng cần thử khảo sát một cách có hệ thống các phương pháp tư
duy, nhằm mở rộng hướng suy nghĩ. Cần chú ý rằng phương pháp tư duy không phải
là phương pháp giải bài. Trọng tâm ở đây là thảo luận quá trình suy nghĩ về
bài toán, chứ không phải tổng kết thuật toán giải bài.

Trước khi bàn cụ thể về các phương pháp tư duy, hãy nhìn lại một phần trong
``bảng cách giải bài toán'' do nhà toán học G. Polya nêu ra năm 1944:
\begin{quote}
Bạn đã từng gặp bài toán này chưa? Bạn có từng gặp cùng một bài toán nhưng với
hình thức hơi khác chưa?

Bạn có biết một bài toán liên quan không? Bạn có biết một định lý nào có thể
dùng được không?

Hãy nhìn vào ẩn số. Hãy thử nghĩ đến một bài toán quen thuộc có cùng ẩn số
hoặc ẩn số tương tự.

Đây là một bài toán liên quan đến bài toán hiện tại của bạn và đã được giải từ
trước.

Bạn có thể dùng nó không? Bạn có thể dùng kết quả của nó không? Bạn có thể
dùng phương pháp của nó không? Để dùng được nó, bạn có cần đưa thêm yếu tố phụ
trợ nào không?

Bạn có thể phát biểu lại bài toán không? Bạn có thể phát biểu lại nó bằng một
cách khác không? Hãy quay về định nghĩa.

Nếu không giải được bài toán đã nêu, trước hết hãy giải một bài toán liên
quan. Bạn có thể nghĩ ra một bài toán liên quan dễ bắt đầu hơn không? Một bài
toán tổng quát hơn? Một bài toán đặc biệt hơn? Một bài toán tương tự? Bạn có
thể giải một phần của bài toán không? Nếu cần, bạn có thể thay đổi ẩn số hoặc
dữ kiện, hoặc thay đổi cả hai, để ẩn số mới và dữ kiện mới gần nhau hơn không?
\end{quote}
Dù bảng này được thiết kế để giải toán, nó đem lại cho chúng ta những gợi mở
rất sâu sắc. Tin học có quan hệ mật thiết với toán học; còn tư duy toán học
lại là một lý thuyết đã trưởng thành, nên nó cung cấp những gợi ý quý báu cho
việc thảo luận của chúng ta.

\section{Các phương pháp tư duy thường gặp}

Phương pháp tư duy rất đa dạng và thay đổi theo từng người. Bài viết này chỉ
chọn vài phương pháp tiêu biểu nhất. Trước hết, theo cách chia truyền thống
``dọc'' và ``ngang'', ta tạm phân loại một số phương pháp tư duy trong thi tin
học như sau:
\begin{itemize}
  \item Theo chiều dọc: thử nghiệm, phỏng đoán và quy nạp; mô hình hóa, hay
  quy về bài toán khác.
  \item Theo chiều ngang: phân loại và chia để trị; tương tự.
\end{itemize}
Trong đó, mô hình hóa và tương tự là hai phương pháp được giới thiệu trọng tâm
ở đây. Tất nhiên, ngoài ra ta cũng ít nhiều chạm tới các phương pháp khác như
đệ quy, sàng lọc, chiến lược tối ưu, tư duy ngược, v.v. Chúng đều quen thuộc
và thường gặp, nhưng vì khuôn khổ có hạn nên sẽ được lược qua.

\subsection{Thử nghiệm, phỏng đoán và quy nạp}

Cấu trúc tri thức, kinh nghiệm giải bài, phẩm chất tư duy và cả tính cách của
mỗi người đều khác nhau. Vì thế, khi đối mặt với một bài toán lạ, trực giác và
``linh cảm'' sinh ra cũng rất muôn hình muôn vẻ. Khi phân tích bài toán, ta
thường không chỉ một lần tự phỏng đoán: ``Liệu có phải thế này không?'',
``Đây chẳng phải là bài nọ hay sao?''. Đó là tư duy phỏng đoán dưới tác động
của tri thức, kinh nghiệm và trí lực cá nhân.

Đôi khi bài toán quá phức tạp, không biết bắt đầu từ đâu và không liên hệ được
với tri thức đã có. Khi ấy ta thường muốn ``làm hết sức mình, trước hết giải
quy mô nhỏ hơn để xem bài toán có quy luật gì không''. Ý nghĩ ``thu nhỏ quy
mô'' và ``tìm quy luật'' chính là tư duy thử nghiệm. Thử nghiệm và phỏng đoán
thường kết hợp với nhau. Không nên cho rằng các lối tư duy ấy là ``không chính
quy'' hay không phải ``chính đạo''; trái lại, chúng thể hiện vẻ đẹp của những
tia sáng tư duy. Quy nạp là quá trình tổng kết hoặc chứng minh các thử nghiệm
và phỏng đoán. Quy nạp thường phải chặt chẽ. Tuy nhiên, trong bối cảnh thi đấu
căng thẳng và gấp gáp, quy nạp chặt chẽ trong lúc thi là việc không đơn giản
và cũng không thường gặp.

Hãy xem một ví dụ về phỏng đoán. Với bài đầu tiên của IOI 1999, ``Little Shop
of Flowers'' (xem Ví dụ~\ref{ex:flowers}), thí sinh giàu kinh nghiệm có thể
chưa đọc xong đề đã đoán rằng phương pháp giải là quy hoạch động. Phỏng đoán
đó dẫn dắt họ nhanh chóng hoàn tất phân tích đề và cuối cùng xác định thuật
toán. Nó bắt nguồn từ nền tảng cơ bản vững chắc, thực hành rộng và kinh nghiệm
phong phú: thí sinh đã luyện nhiều bài quy hoạch động, nên khi suy nghĩ có thể
tự nhiên liên hệ tương tự với các ví dụ đã làm. Vì vậy phỏng đoán tuyệt đối
không phải là đoán mò.

Một loại phỏng đoán khác bắt nguồn từ thử nghiệm. Dưới đây là một ví dụ rất
kinh điển.

\paragraph{Ví dụ 1.}
Cho \(m,n\in\mathbb{N}\), \(1\le m,n\le k\le 10^9\), và
\[
  (n^2-mn-m^2)^2=1.
\]
Nhập \(k\), hãy tìm \(m,n\) sao cho \(m^2+n^2\) lớn nhất. (NOI 1995)

Từ quy mô dữ liệu, ta có thể ``đoán'' rằng bài này chắc chắn ẩn chứa một quan
hệ toán học đơn giản hơn. Nhưng quan hệ toán học trong đề không lộ rõ, và phân
tích toán học trực tiếp khá khó. Trong khi đó, chỉ bằng hai vòng lặp đơn giản,
ta có thể dễ dàng liệt kê các trường hợp nhỏ:
\[
\begin{array}{c|cccccc}
k & 1 & 2 & 3,4 & 5,6,7 & 8,9,10,11,12 & 13,\ldots\\
\hline
m & 1 & 2 & 3 & 5 & 8 & 13,\ldots\\
n & 1 & 1 & 2 & 3 & 5 & 8,\ldots
\end{array}
\]
Qua các thử nghiệm này, ta phỏng đoán rằng các cặp \(m,n\) thỏa điều kiện tạo
thành dãy Fibonacci. Thật vậy:
\[
\begin{aligned}
n^2-mn-m^2
  &=-(m^2+mn-n^2)\\
  &=-\bigl((m+n)^2-mn-2n^2\bigr)\\
  &=-\bigl((m+n)^2-n(m+n)-n^2\bigr).
\end{aligned}
\]
Nhiều phỏng đoán khác đến từ phép tương tự. Tóm lại, phỏng đoán là một dạng
suy luận hợp lý, có lợi cho việc dẫn dắt hướng suy nghĩ~\cite{math-thinking}.

\subsection{Mô hình hóa}

``Bạn có thể phát biểu lại bài toán không? Bạn có thể phát biểu lại nó bằng
một cách khác không?''

Thế giới khách quan vốn phức tạp. Khi đối mặt với một bài toán mới, suy nghĩ
thông thường là thông qua phân tích, liên tục biến đổi và chuyển hóa để thu
được một bài toán có bản chất giống hệt nhưng quen thuộc hơn, hoặc trừu tượng
và tổng quát hơn. Đó là tư tưởng quy về bài toán khác. Ta gọi bài toán hoặc
đối tượng ban đầu là nguyên mẫu, còn đối tượng hoặc bài toán tương đối định
hình, được mô phỏng hoặc lý tưởng hóa sau khi quy đổi là mô hình. Mô hình và
nguyên mẫu phải có bản chất giống nhau, tức là tương đương. Các hướng mô hình
hóa chủ yếu gồm mô hình đồ thị, mô hình toán học và mô hình quy hoạch, đặc
biệt là quy hoạch động.

\subsubsection{Mô hình đồ thị}

Mô hình hóa bằng đồ thị là tư tưởng ta dùng thường xuyên nhất. Lý thuyết đồ
thị có rất nhiều thuật toán đã trưởng thành. Chuyển một bài toán lạ thành bài
toán đồ thị trừu tượng cũng là việc nhiều thí sinh ưa làm. Lý do là sau khi
đồ thị hóa, quan hệ dữ liệu vốn phức tạp và rối rắm trở nên ngắn gọn, sáng
sủa, và hướng giải bài cũng rõ ràng hơn.

\paragraph{Ví dụ 2. Bài ``Optimal Connected Subset''.}
Trong tập các điểm nguyên trên mặt phẳng tọa độ, hai điểm cách nhau \(1\) được
gọi là kề nhau. Nếu \(S\) là một tập điểm nguyên, \(R,T\) là các điểm thuộc
\(S\), và trong \(S\) tồn tại một dãy điểm đôi một khác nhau
\[
  R-Q_1-Q_2-\cdots-Q_k-T
\]
sao cho hai điểm liên tiếp bất kỳ đều kề nhau, thì gọi
\(Q_1-Q_2-\cdots-Q_k\) là một đường trong tập điểm nguyên \(S\) nối \(R\) với
\(T\). Nếu giữa hai điểm bất kỳ của \(S\) đều có đúng một đường nối, thì \(S\)
được gọi là tập điểm nguyên đơn. Cho \(S\) và trọng số của mọi điểm trong
\(S\), hãy tìm tập điểm nguyên đơn \(B\), là tập con của \(S\), sao cho tổng
trọng số của \(B\) là lớn nhất. (NOI 1999)

Bản thân phát biểu bài này đã rất phức tạp, như thể cố ý dựng nên một đống
quan hệ rắc rối. Vì vậy ta tự nhiên muốn ``đơn giản hóa'' nó. Rất dễ nghĩ tới
đồ thị. Nhưng ``đồ thị'' thông thường không có tính chất đặc biệt của ``tập
điểm nguyên đơn'' trong đề: giữa hai điểm bất kỳ có đúng một đường nối. Hãy
chú ý rằng tính chất này đúng bằng ``trong đồ thị không có chu trình'', mà cây
chính là đồ thị không có chu trình. Khi tư tưởng xây dựng mô hình ``cây'' được
xác lập, bài toán đã được đơn giản hóa. Theo đó là hàng loạt thuật toán liên
quan đến cây như duyệt cây, đếm trên cây, v.v. Nhờ vậy hướng giải của cả bài
rõ hơn rất nhiều. Một hướng là: ``tập con liên thông tối ưu'' nhất định là một
cây con của một cây tìm kiếm sâu, với gốc là một điểm nào đó trong ``tập điểm
nguyên đơn'' ban đầu. Do đó, chỉ cần duyệt ``tập điểm nguyên đơn'' ban đầu một
lần, đồng thời tính tổng trọng số lớn nhất của mỗi cây con.

\paragraph{Ví dụ 3. Bài ``Mèo bắt chuột''.}
Cho một chiếc lồng, trong đó ký hiệu \(\ast\) là chướng ngại vật. Vị trí ban
đầu của mèo đã biết, quỹ đạo di chuyển của chuột đã biết. Mỗi giây chuột đi
một bước, mèo đi \(S\) bước. Hỏi mèo có bắt được chuột không? (GDOI 1997)

Muốn bắt được chuột, mèo phải đến một vị trí nào đó trên quỹ đạo của chuột
trước chuột. Vì vậy cần tính thời gian ngắn nhất để mèo đến mọi vị trí trên
quỹ đạo của chuột. Ta có thể liên tưởng đến phương pháp gán nhãn trong đường
đi ngắn nhất, đây lại là một phỏng đoán. Đó là xét từ phía ``chưa biết'' của
bài toán. Xét tiếp từ phía ``đã biết'': mèo đi trong lồng, độ dài bước đã biết,
thế nên khi vị trí của nó xác định thì các ô nó có thể đến ở bước tiếp theo
cũng xác định. Như vậy, coi mỗi ô là một đỉnh, nối cạnh giữa hai ô nếu có thể
đi từ ô này sang ô kia trong một bước, ta thu được một đồ thị. Dễ thấy đồ thị
này tương đương với bài toán gốc. Đến đây mô hình hóa đã hoàn tất và phác
thảo lời giải đã hình thành.

Trong quá trình tư duy này không chỉ có mô hình hóa, mà còn mang màu sắc của
phỏng đoán và tư duy hai chiều.

\paragraph{Ví dụ 4. Bài ``Patches vs. Bugs''.}
Một phần mềm có \(n\) lỗi (\(1\le n\le 20\)) và cần một số bản vá. Mỗi bản vá
chỉ dùng được khi phần mềm chứa một số lỗi nhất định và không chứa một số lỗi
khác; mỗi bản vá chỉ sửa một phần lỗi, đồng thời có thể gây ra một số lỗi
khác. Thời gian chạy của mỗi bản vá khác nhau. Ta muốn sửa hết lỗi của phần
mềm với tổng thời gian nhỏ nhất. (Tuyển chọn đội tuyển Trung Quốc IOI 1999)

Vì bài toán có yêu cầu ``ít thời gian nhất'', ta tự nhiên nghĩ tới đường đi
ngắn nhất. Để đạt mục tiêu ấy, ta thử mô hình hóa bài toán. Trạng thái của
\(n\) lỗi, tức có hay không có từng lỗi, có thể biểu diễn ngắn gọn bằng một số
nhị phân; tất cả các trạng thái có thể có tạo thành các đỉnh của đồ thị. Cạnh
có hướng giữa các đỉnh biểu thị việc có thể dùng một bản vá để chuyển sang
trạng thái khác. Trọng số cạnh là thời gian chạy của bản vá. Khi đó bài toán
được trừu tượng hóa thành tìm đường đi ngắn nhất từ nguồn \(2^n-1\) đến đích
\(0\). Mô hình này hoàn toàn tương đương với nguyên mẫu và thể hiện đúng
``bản chất'' của nguyên mẫu, nên có thể xem là một mô hình rất đẹp.

Nói chung, với những bài liên quan đến đại lượng rời rạc và quan hệ dữ liệu
đan xen phức tạp, ta có thể cân nhắc mô hình hóa bằng đồ thị. Nếu đồ thị hóa
thành công, ta dễ ước lượng bậc phức tạp đại khái của thuật toán: đa thức,
mũ, hay giai thừa.

Khi xây dựng mô hình đồ thị cần chú ý: mô hình càng gần bản chất của nguyên
mẫu càng tốt. Nếu mô hình đặc biệt hơn nguyên mẫu, ta có thể mất nghiệm hoặc
không tìm được nghiệm. Nếu mô hình tổng quát hơn nguyên mẫu, thuật toán thu
được có thể có độ phức tạp quá cao, dễ quá thời gian và cũng không giải được.
Hãy xem ví dụ sau.

\paragraph{Ví dụ 5. Bài ``Hidden Codes''.}
Cho một từ điển mã \texttt{words.inp} và một văn bản \texttt{text.inp}. Văn bản
chứa một số mã; các chữ cái của mỗi mã xuất hiện trong văn bản theo đúng thứ
tự nhưng không nhất thiết liên tiếp. Cần nhận diện các mã này trong văn bản.
Các cách nhận diện khác nhau có thể thu được các mã khác nhau. Yêu cầu các
dãy con trong văn bản chứa mã, gọi là dãy phủ, không giao nhau và không chồng
lên nhau, đồng thời tổng độ dài của tất cả các mã là lớn nhất. Ví dụ có bốn
mã \texttt{RuN}, \texttt{RaBbit}, \texttt{HoBbit}, \texttt{StoP}, và văn bản
\texttt{StXRuYNvRuHoaBbvizXztNwRRuuNNP}; khi đó có thể chọn ba dãy phủ chứa
ba mã \texttt{RuN}, \texttt{RaBbit}, \texttt{RuN}, với tổng độ dài bằng \(12\).
(IOI 1999)

Bài này yêu cầu tìm một nhóm dãy phủ chứa lượng mã nhiều nhất. Qua phân tích,
ta biết rằng các dãy phủ ``tối tiểu bên phải'', nghĩa là không tiền tố nào của
nó đã là dãy phủ, đều xác định được trong văn bản. Vì thế ta cân nhắc lấy mỗi
dãy phủ ``tối tiểu bên phải'' làm một đỉnh; cạnh giữa hai đỉnh biểu thị hai
dãy không chồng lên nhau; độ dài mã mà mỗi dãy chứa ứng với giá trị của đỉnh.
Bài toán chuyển thành tìm chuỗi dài nhất trong đồ thị.

Rõ ràng đồ thị hóa như vậy không ổn: ta đều biết bài toán đường đi dài nhất
không có thuật toán tốt. Vì sao sau khi đồ thị hóa, độ phức tạp của bài toán
lại tăng cao như vậy? Nguyên nhân là trong nguyên mẫu, các dãy phủ có thứ tự,
tức được sắp xếp tuyến tính. Còn đồ thị thì ``không gian'' hơn; giữa các đỉnh
của đồ thị không có quan hệ trước sau tự nhiên. Biến một cấu trúc ``tuyến
tính'' thành một cấu trúc ``không gian'' không những không tương đương với
nguyên mẫu, mà còn tổng quát hóa quá mức, tất yếu làm độ phức tạp của bài toán
tăng thêm một bậc.

Dù vậy, không phải mọi mô hình hóa đều phải tương đương chặt chẽ. Mô hình đôi
khi chỉ là một phép tương tự của nguyên mẫu; thông thường ta vẫn không thể
tách khỏi nguyên mẫu khi thiết kế thuật toán cụ thể. Chẳng hạn trong
Ví dụ~2, ``tập điểm nguyên đơn'' thực ra còn đặc biệt hơn cây: gốc có nhiều
nhất \(4\) con, các đỉnh còn lại có nhiều nhất \(3\) con, và giữa các đỉnh còn
có những quan hệ phức tạp hơn. Khi bỏ qua các hạn chế đó, mô hình cây thu được
tất nhiên không tương đương chặt chẽ với nguyên mẫu; nó tổng quát hơn nguyên
mẫu. Nhưng mô hình ấy và nguyên mẫu giống nhau đến mức bản chất của chúng nhất
quán. Trái lại, nếu hiểu ``tập điểm nguyên đơn'' là một đồ thị tùy ý thì mô
hình quá tổng quát và mất tác dụng. Tương tự, trong Ví dụ~3, không phải đồ thị
bất kỳ nào cũng có thể chuyển thành cấu trúc ``lồng'', tức lồng và đồ thị
không tương đương chặt chẽ; nhưng điều đó không cản trở việc giải bài toán.

\subsubsection{Mô hình toán học}

Trong các mô hình toán học, được dùng nhiều nhất là các mô hình tổ hợp. Tổ hợp
là lĩnh vực rất khó; xây dựng mô hình tổ hợp đòi hỏi nền tảng toán học tốt và
thể hiện rõ nhất sự chặt chẽ của tư duy. Đồng thời nó cũng liên quan đến phân
loại, chia để trị, đệ quy và các phương pháp tư duy khác.

Các loại bài toán đếm đều có thể cân nhắc xây dựng mô hình tổ hợp. Mô hình tổ
hợp thường dùng nhất là quan hệ đệ quy. Dùng trực tiếp quan hệ đệ quy để đếm
là khả thi, nhưng nếu giải quan hệ đệ quy trước rồi mới đếm thì có thể giảm
độ phức tạp rất nhiều.

Ví dụ~1 dùng dãy Fibonacci trong tổ hợp làm mô hình. Mô hình tổ hợp kinh điển
nhất có lẽ là số Catalan. Một nguyên mẫu của số Catalan như sau.

\paragraph{Ví dụ 6.}
Có tích liên tiếp của \(n\) số \(k_1k_2k_3\cdots k_n\). Hãy chèn đủ dấu ngoặc
sao cho mỗi tích con đúng bằng tích của hai thừa số. Ví dụ, tích
\(k_1k_2k_3k_4\) có thể được đặt ngoặc thành
\[
((k_1k_2)(k_3k_4)),\quad (((k_1k_2)k_3)k_4),\quad
((k_1(k_2k_3))k_4),\quad (k_1((k_2k_3)k_4)),\quad
(k_1(k_2(k_3k_4))).
\]
Chú ý rằng không được thay đổi thứ tự của các \(k_i\). Hãy tính số cách
\(a_n\) để đặt ngoặc cho tích \(k_1k_2k_3\cdots k_n\).

Khi xây dựng quan hệ đệ quy, thường phải tìm cách tách bài toán gốc thành các
bài toán con có quy mô nhỏ hơn và có cùng tính chất. Dễ biết \(a_2=1\),
\(a_3=2\), ứng với \(((k_1k_2)k_3)\) và \((k_1(k_2k_3))\). Đặt \(a_0=0\),
\(a_1=1\). Xét phép nhân cuối cùng của tích \(n\) số, tức cặp ngoặc ngoài
cùng. Phép nhân cuối cùng này liên quan đến tích của hai bài toán con:
\[
  ((k_1k_2\cdots k_i)(k_{i+1}k_{i+2}\cdots k_n)).
\]
Số cách đặt ngoặc cho hai bài toán con lần lượt là \(a_i\) và \(a_{n-i}\), nên
hai bài toán con có tổng cộng \(a_i a_{n-i}\) cách, trong đó \(i\) chạy từ
\(1\) đến \(n-1\). Cộng theo mọi \(i\), ta nhận được quan hệ đệ quy, với
\(n\ge 2\):
\[
  a_n=a_1a_{n-1}+a_2a_{n-2}+\cdots+a_ia_{n-i}+\cdots+a_{n-1}a_1.
  \tag{*}
\]

Đây chính là quan hệ đệ quy của số Catalan. Giải quan hệ này ta được
\[
  a_n=\frac{1}{n}C_{2n-2}^{\,n-1}.
\]
Có rất nhiều bài toán đếm có thể quy về quan hệ đệ quy \((*)\) của số
Catalan, tức số Catalan có nhiều nguyên mẫu phong phú. Ví dụ:
\begin{itemize}
  \item \term{Bài toán đa giác lồi.} Chia một đa giác lồi \(n\) cạnh bằng
  \(n-3\) đường chéo không cắt nhau thành \(n-2\) tam giác. Hỏi có bao nhiêu
  cách chia?
  \item \term{Bài toán đếm cây.} Có bao nhiêu cây nhị phân với \(n\) nút?
\end{itemize}
Dùng chia để trị, đệ quy và các phương pháp tương tự, ta dễ thu được cùng quan
hệ đệ quy \((*)\) cho hai ví dụ này.

Bài toán tổ hợp biến hóa muôn hình, phương pháp giải cũng đa dạng. Mô hình hóa
chỉ là một hướng trong đó, đôi khi chỉ là một bước trong quá trình giải, nhưng
qua các ví dụ trên có thể thấy tác dụng của nó rất lớn.

\subsubsection{Mô hình quy hoạch}

Mô hình quy hoạch, tức quy hoạch toán học, là một loại mô hình toán học. Vì nó
thường gặp nên được tách ra trình bày riêng.

Các mô hình quy hoạch quen thuộc chủ yếu gồm quy hoạch nguyên và quy hoạch
động. Quy hoạch nguyên là bài toán quy hoạch mà mọi biến đều nhận giá trị
nguyên; thuật toán cho loại bài này thường có bậc giai thừa. Biến trạng thái
quyết định của quy hoạch động cũng nhận giá trị nguyên, nhưng độ phức tạp thuật
toán lại có thể là đa thức.

Với bài toán giải nhiều biến có ràng buộc, đặc biệt khi điều kiện ràng buộc
đòi một hàm đạt giá trị lớn nhất hoặc nhỏ nhất, ta có thể cân nhắc xây dựng mô
hình quy hoạch. Khi xây dựng mô hình này, quy hoạch động là lựa chọn đầu tiên;
nếu chọn quy hoạch nguyên thì cần chú ý tối ưu thuật toán.

\paragraph{Ví dụ 7. Bài ``Little Shop of Flowers''.}
\label{ex:flowers}
Có \(F\) bó hoa cần cắm từ trái sang phải vào \(V\) bình
\((1\le F\le V\le 100)\). Hoa và bình lần lượt được đánh số bằng các số nguyên
\(1\) đến \(F\) và \(1\) đến \(V\). Hoa phải được sắp theo thứ tự tăng dần của
chỉ số hoa. Mỗi cách cắm một bông hoa vào một bình khác nhau sinh ra một giá
trị thẩm mỹ khác nhau, trong khoảng từ \(-50\) đến \(50\). Hãy tìm giá trị
thẩm mỹ lớn nhất và một cách sắp xếp đạt giá trị đó. (IOI 1999)

Thực chất đây là việc giải một dãy biến nguyên, và hàm mục tiêu thuộc dạng
``lớn nhất''. Vì vậy có thể xây dựng mô hình quy hoạch:
\[
\begin{gathered}
  x_i\in\mathbb{N},\\
  x_i<x_{i+1}\quad (1\le i<F),\\
  \max f=\sum_i A_{i,x_i}.
\end{gathered}
\]
Đây là quy hoạch nguyên. Nếu quyết định các \(x_i\) theo từng giai đoạn, mô
hình sẽ chuyển dần thành quy hoạch động.

\paragraph{Ví dụ 8. Bài ``Chia bàn cờ''.}
Từ một bàn cờ hình chữ nhật, cắt một nhát để lấy ra một bàn cờ hình chữ nhật
và phần còn lại cũng là hình chữ nhật. Tiếp tục cắt phần còn lại như vậy. Sau
\(n-1\) lần cắt, thu được \(n\) bàn cờ hình chữ nhật. Mỗi ô trên bàn cờ có một
điểm số. Cần chia bàn cờ theo quy tắc trên thành \(n\) bàn cờ hình chữ nhật
sao cho phương sai bình phương trung bình của các phần là nhỏ nhất. (NOI 1999)

Nhiều thí sinh chọn quy hoạch động hoặc thuật toán tìm kiếm. Vì sao? Bởi tuy
đề chỉ yêu cầu một giá trị tối ưu, ta nên nhận ra rằng phương sai được quyết
định bởi cách cắt, tức bởi \(n-1\) biến đại diện cho vị trí của từng ``nhát
cắt''. Vì vậy bài toán ẩn chứa việc giải \(n-1\) biến này. So với các mô hình
khác, bài toán giống mô hình quy hoạch hơn. Cốt lõi của lời giải chính là giá
trị của \(n-1\) biến ẩn đó. Nói chung, tìm kiếm luôn là hướng dễ nghĩ tới, tất
nhiên vẫn cần chú ý tối ưu. Còn nếu thêm được khái niệm giai đoạn và chủ động
chia việc giải \(n-1\) biến này thành \(n-1\) giai đoạn, ta sẽ tự nhiên tiến
gần đến mô hình quy hoạch động.

\subsection{Phân loại và chia để trị}

Phân loại và chia để trị bắt nguồn từ mong muốn biến việc lớn thành việc nhỏ,
rồi xử lý việc nhỏ.

Phân loại là chia các phần của bài toán theo logic. Mô-đun hóa chương trình là
một sản phẩm của tư tưởng phân loại; thiết kế dữ liệu kiểm thử cũng thể hiện
tư tưởng phân loại. Quá trình tư duy phân loại phải toàn diện và chặt chẽ. Khi
biến đổi và chuyển hóa bài toán, cần bảo đảm đã xét mọi mặt của bài toán,
không lặp và không sót. Không chỉ bài toán tổ hợp, mà mọi loại bài toán khác
cũng cần chú ý điều này.

\paragraph{Ví dụ 9.}
Trên mặt phẳng có \(n\) điểm (\(n>1\)). Cần tìm hình tròn có diện tích nhỏ nhất
có thể phủ tất cả các điểm đó.

Một hướng suy nghĩ là: vì ba điểm xác định một đường tròn, khi đường tròn phủ
tất cả các điểm co nhỏ dần thì chắc chắn có xu hướng ``co'' đến mức một số
điểm vừa nằm trên đường tròn, có thể liên tưởng đến bao lồi. Vì vậy để xác
định đường tròn này, chỉ cần liệt kê ba điểm, độ phức tạp \(O(n^3)\).

Hướng suy nghĩ này có khiếm khuyết. Xu hướng ``co'' quả thực tồn tại; nếu cuối
cùng có không dưới ba điểm nằm trên đường tròn thì đương nhiên có thể liệt kê.
Vấn đề là: có thể nào cuối cùng chỉ có hai điểm, hoặc chỉ một điểm, nằm trên
đường tròn hay không? Nếu chỉ có một điểm thì rõ ràng đường tròn còn có thể
tiếp tục ``co''. Nếu chỉ có hai điểm, quá trình ``co'' sẽ khiến tâm đường tròn
nằm ở trung điểm đoạn nối hai điểm đó, tức đoạn thẳng nối hai điểm là đường
kính của đường tròn. Lúc này hoàn toàn có thể không có điểm nào khác nằm trên
đường tròn; nếu dùng ba điểm để xác định đường tròn thì ngược lại sẽ làm nó
lớn hơn.

Chia để trị đã nhiều lần được nhắc tới ở trên. Chia để trị là một dạng đặc
biệt của phân loại, thường gắn với đệ quy. Đặc điểm của nó là các bài toán con
sau khi chia vẫn là bài toán gốc với quy mô nhỏ hơn, nên mạch suy nghĩ rất rõ
ràng và ngắn gọn. Có thể nói tư tưởng chia để trị đã hòa vào nhiều tư tưởng
khác ở trên, như tổ hợp, quy hoạch động và tìm kiếm.

\subsection{Tương tự}

``Bạn có từng gặp cùng một bài toán nhưng với hình thức hơi khác chưa? Bạn có
biết một bài toán liên quan không? Bạn có biết một định lý nào có thể dùng
được không?'' ``Hãy thử nghĩ đến một bài toán quen thuộc có cùng ẩn số hoặc ẩn
số tương tự.'' ``Đây là một bài toán liên quan đến bài toán hiện tại của bạn
và đã được giải từ trước. Bạn có thể dùng nó không?'' ``Một bài toán tương tự
thì sao?''

Chỉ cần nhìn trong bảng của Polya có nhiều câu liên quan đến tương tự như vậy
cũng đủ thấy tầm quan trọng của nó. Nếu những phần trước còn mang màu sắc phân
tích phương pháp giải bài, thì tương tự là một phương pháp tư duy thuần túy.
Thực ra, việc bài viết này ngay từ đầu mượn bảng giải toán của Polya đã là một
phép tương tự xuyên ngành; trong các phần trước, từ ``tương tự'' cũng đã nhiều
lần xuất hiện.

Trong Ví dụ~2, vì sao nghĩ tới cây? Bởi ta đem tính chất ``giữa hai điểm trong
tập điểm nguyên đơn có đúng một đường'' trong đề so sánh với tính chất của
cây: giữa hai đỉnh bất kỳ của cây có đúng một đường.

Trong Ví dụ~3 và Ví dụ~4, vì bài toán liên quan đến ``ngắn nhất'', ta tương tự
với đường đi ngắn nhất, tìm chỗ giống nhau, liên tục chuyển hóa bài toán về
mục tiêu, và cuối cùng hoàn thành mô hình hóa. Vấn đề tương tự cũng xuất hiện
trong bài ``Traffic Lights'' của IOI 1999. Bài đó cũng yêu cầu ``ngắn nhất'',
nên cũng tương tự với đường đi ngắn nhất, tìm khoảng cách giữa hai bài toán,
rồi thử loại bỏ khoảng cách đó và chuyển hóa về mục tiêu; kết quả là nó cũng
thỏa tính không hậu hiệu.

Tương tự có thể là tương tự bản chất, cũng có thể là tương tự hình thức; miễn
là gợi mở được hướng suy nghĩ thì đều có giá trị. Ví dụ~5 hơi giống bài
``Alien'' của IOI 1998, nên có thể nghĩ đến việc mượn một phần tư tưởng của
bài đó, như cách lưu trữ, đọc dữ liệu ``một lần'', hoặc tư tưởng quy hoạch
động. Không phải ý nào cũng nhất định hữu ích, nhưng mức độ gợi mở là đáng kể.

Ví dụ~6 có thể xem là ví dụ điển hình. Mỗi khi làm bài tổ hợp, nếu nhớ đến bài
này, nhớ quan hệ truy hồi được xây dựng ra sao và được giải như thế nào, ta có
thể tìm thấy khá nhiều điểm để tham khảo.

Hình thức của Ví dụ~8 cũng rất thường gặp, chẳng hạn: ``tìm một ma trận con
của một ma trận số nguyên sao cho tổng các phần tử của ma trận con là lớn
nhất''. Từ đó cũng có thể gợi ra tư tưởng chia bài toán thành các bài toán con.

Dĩ nhiên, vì tương tự là một lối tư duy không chặt chẽ, tương tự thất bại là
điều khó tránh. Chẳng hạn trong Ví dụ~5, phép tương tự với mô hình đồ thị quen
thuộc là không thành công. Điều này một lần nữa chứng tỏ tương tự chỉ là công
cụ gợi mở hướng suy nghĩ.

Cần chú ý rằng tương tự phải có ``vật để so''. Nghĩa là trước hết phải coi
trọng tích lũy tri thức và thực hành. Tri thức càng phong phú, thực hành càng
rộng thì càng có nhiều đối tượng để so sánh tương tự; hướng suy nghĩ càng rộng
và sự gợi mở càng lớn.

Tóm lại, tương tự là một dạng của sự giống nhau. Vì ý nghĩa của ``giống nhau''
rất rộng, hình thức biểu hiện của phương pháp tương tự cũng phong phú và đa
dạng. Tương tự cũng là một phương pháp phát hiện để tìm hướng giải bài, phỏng
đoán đáp án hoặc kết luận của bài toán~\cite{math-thinking}. Dù kết quả của
tương tự không nhất định đúng, nó là tư duy phân kỳ, là nền tảng của tư duy
sáng tạo và có vai trò rất quan trọng.

\section{Kết luận}

Khi đối mặt với bài toán, thông thường ta luôn thông qua quan sát để làm rõ
bài toán, nắm lấy đặc trưng của đề, rồi liên tưởng, tìm kiếm và hồi tưởng trên
diện rộng. Nói cách khác, dựa vào tri thức và kinh nghiệm đã có, ta đưa ra
những hiểu biết và phán đoán mang tính trực giác, chọn hướng suy nghĩ tổng thể
hoặc điểm bắt đầu. Có tìm được chiến lược phù hợp hay không liên quan đến góc
độ quan sát bài toán cũng như độ rộng, hẹp, sâu, nông của phạm vi liên tưởng.
Khi tư duy bị cản trở, nên điều chỉnh hướng suy nghĩ, đổi góc nhìn rồi tiếp
tục phân tích cho đến khi nảy sinh hướng mới~\cite{math-thinking}. Trong quá
trình đó, các phương pháp tư duy không được dùng một cách cô lập, mà được vận
dụng một cách biện chứng.

Cuối cùng, mục đích của bài viết khi thảo luận phương pháp tư duy có hai mặt:
một là gợi mở hướng suy nghĩ cho người đọc; hai là chỉ ra rằng tích lũy tri
thức và thực hành thường ngày mới là con đường tốt nhất để phát triển tư duy
cá nhân.

\section*{Phụ lục}

\begin{enumerate}
  \item Bảng của Polya gồm bốn phần: hiểu rõ bài toán, lập kế hoạch, thực hiện
  kế hoạch và nhìn lại. Ở đây chỉ trích một phần của mục lập kế hoạch.
  \item Có thể xem thêm ví dụ trong bài ``Bài toán hệ số đệ quy của hàm lũy
  thừa'' của tỉnh Phúc Kiến, đăng trên tạp chí \textit{Olympic Tin học}, số 3
  năm 1999, trang 26.
  \item Thực ra Ví dụ~4 còn cần xử lý vấn đề dung lượng lưu trữ
  (\(2^{20}=1{,}048{,}576\)), liên quan đến kỹ thuật lưu trữ trên máy tính.
  Điều này không nằm trong phạm vi thảo luận của bài viết nên được lược qua.
  \item Việc giải quan hệ đệ quy thuộc tổ hợp, không nằm trong phạm vi bài
  viết, nên được lược qua. Ở đây chỉ bàn cách thiết lập quan hệ đệ quy. Cách
  giải số Catalan có trong nhiều sách tổ hợp, xin tham khảo riêng.
  \item Xem \textit{Lý luận tư duy toán học}, trang 141.
  \item Xem \textit{Lý luận tư duy toán học}, trang 242.
\end{enumerate}

\section*{Tài liệu tham khảo}

\begin{thebibliography}{9}
\bibitem{math-thinking}
Ren Zhanghui.
\textit{Lý luận tư duy toán học}.
Nhà xuất bản Giáo dục Quảng Tây, ấn bản 1996.

\bibitem{informatics-olympiad}
Tạp chí \textit{Olympic Tin học}.

\bibitem{comb-book}
Zhuang Xingu.
\textit{Tổ hợp toán học và ứng dụng trong khoa học máy tính}.
Nhà xuất bản Đại học Khoa học Kỹ thuật Điện tử Tây An, ấn bản 1989.
\end{thebibliography}

\end{document}
